L'algorithmique du collège tient en trois idées : une variable prend une valeur, une boucle
répète, un test choisit. Ce chapitre les traite comme des mathématiques, c'est-à-dire qu'il
démontre ce que les programmes calculent.

Deux programmes de calcul qui donnent le même résultat pour toute entrée, une boucle dont on
établit l'invariant, un algorithme dont on prouve qu'il termine : ce sont là des énoncés, et
ils se démontrent. En revanche \og{} écrire un programme qui \dots{} \fg{} est une tâche.
La différence est la même qu'entre \og{} tracer la médiatrice \fg{} et \og{} démontrer sa
caractérisation \fg{}, et elle traverse tout ce livre.
